\documentclass[11pt,a4paper]{article}
% preamble.tex — Shared LaTeX preamble for all tcfs design documents
% Usage: % preamble.tex — Shared LaTeX preamble for all tcfs design documents
% Usage: % preamble.tex — Shared LaTeX preamble for all tcfs design documents
% Usage: \input{preamble} at top of each .tex file

% Encoding and fonts
\usepackage[T1]{fontenc}
\usepackage[utf8]{inputenc}
\usepackage{lmodern}
\usepackage{textcomp}

% Page geometry
\usepackage[margin=2.5cm]{geometry}

% Colors and hyperlinks
\usepackage[dvipsnames]{xcolor}
\usepackage[colorlinks=true,linkcolor=NavyBlue,urlcolor=NavyBlue,citecolor=NavyBlue]{hyperref}

% Code listings
\usepackage{listings}
\lstset{
  basicstyle=\ttfamily\small,
  breaklines=true,
  frame=single,
  backgroundcolor=\color{gray!10},
  xleftmargin=2em,
  framexleftmargin=1.5em,
  numbers=none,
  showstringspaces=false,
  columns=flexible,
  literate={←}{$\leftarrow$}1 {→}{$\rightarrow$}1 {⟶}{$\longrightarrow$}1,
}
\lstdefinelanguage{json}{
  string=[s]{"}{"},
  comment=[l]{//},
  morecomment=[s]{/*}{*/},
  literate=
    *{:}{{{\color{NavyBlue}:}}}{1}
    {,}{{{\color{NavyBlue},}}}{1}
    {\{}{{{\color{NavyBlue}\{}}}{1}
    {\}}{{{\color{NavyBlue}\}}}}{1}
    {[}{{{\color{NavyBlue}[}}}{1}
    {]}{{{\color{NavyBlue}]}}}{1},
}
\lstdefinelanguage{toml}{
  comment=[l]{\#},
  string=[b]",
  morestring=[b]',
}
\lstdefinelanguage{protobuf}{
  morekeywords={service,rpc,returns,message,string,uint32,uint64,int64,bool,bytes,repeated,optional,oneof,enum,map},
  morecomment=[l]{//},
  morestring=[b]",
}

% Tables
\usepackage{longtable}
\usepackage{booktabs}
\usepackage{array}
\usepackage{tabularx}

% Graphics
\usepackage{graphicx}
\usepackage{float}

% Better lists
\usepackage{enumitem}

% Header/footer
\usepackage{fancyhdr}
\pagestyle{fancy}
\fancyhf{}
\fancyfoot[C]{\thepage}
\renewcommand{\headrulewidth}{0.4pt}

% Tighter spacing
\setlength{\parindent}{0pt}
\setlength{\parskip}{6pt}

% Convenience macros
\newcommand{\tcfs}{\texttt{tcfs}}
\newcommand{\tcfsd}{\texttt{tcfsd}}
\newcommand{\code}[1]{\texttt{#1}}
\newcommand{\filepath}[1]{\texttt{#1}}
 at top of each .tex file

% Encoding and fonts
\usepackage[T1]{fontenc}
\usepackage[utf8]{inputenc}
\usepackage{lmodern}
\usepackage{textcomp}

% Page geometry
\usepackage[margin=2.5cm]{geometry}

% Colors and hyperlinks
\usepackage[dvipsnames]{xcolor}
\usepackage[colorlinks=true,linkcolor=NavyBlue,urlcolor=NavyBlue,citecolor=NavyBlue]{hyperref}

% Code listings
\usepackage{listings}
\lstset{
  basicstyle=\ttfamily\small,
  breaklines=true,
  frame=single,
  backgroundcolor=\color{gray!10},
  xleftmargin=2em,
  framexleftmargin=1.5em,
  numbers=none,
  showstringspaces=false,
  columns=flexible,
  literate={←}{$\leftarrow$}1 {→}{$\rightarrow$}1 {⟶}{$\longrightarrow$}1,
}
\lstdefinelanguage{json}{
  string=[s]{"}{"},
  comment=[l]{//},
  morecomment=[s]{/*}{*/},
  literate=
    *{:}{{{\color{NavyBlue}:}}}{1}
    {,}{{{\color{NavyBlue},}}}{1}
    {\{}{{{\color{NavyBlue}\{}}}{1}
    {\}}{{{\color{NavyBlue}\}}}}{1}
    {[}{{{\color{NavyBlue}[}}}{1}
    {]}{{{\color{NavyBlue}]}}}{1},
}
\lstdefinelanguage{toml}{
  comment=[l]{\#},
  string=[b]",
  morestring=[b]',
}
\lstdefinelanguage{protobuf}{
  morekeywords={service,rpc,returns,message,string,uint32,uint64,int64,bool,bytes,repeated,optional,oneof,enum,map},
  morecomment=[l]{//},
  morestring=[b]",
}

% Tables
\usepackage{longtable}
\usepackage{booktabs}
\usepackage{array}
\usepackage{tabularx}

% Graphics
\usepackage{graphicx}
\usepackage{float}

% Better lists
\usepackage{enumitem}

% Header/footer
\usepackage{fancyhdr}
\pagestyle{fancy}
\fancyhf{}
\fancyfoot[C]{\thepage}
\renewcommand{\headrulewidth}{0.4pt}

% Tighter spacing
\setlength{\parindent}{0pt}
\setlength{\parskip}{6pt}

% Convenience macros
\newcommand{\tcfs}{\texttt{tcfs}}
\newcommand{\tcfsd}{\texttt{tcfsd}}
\newcommand{\code}[1]{\texttt{#1}}
\newcommand{\filepath}[1]{\texttt{#1}}
 at top of each .tex file

% Encoding and fonts
\usepackage[T1]{fontenc}
\usepackage[utf8]{inputenc}
\usepackage{lmodern}
\usepackage{textcomp}

% Page geometry
\usepackage[margin=2.5cm]{geometry}

% Colors and hyperlinks
\usepackage[dvipsnames]{xcolor}
\usepackage[colorlinks=true,linkcolor=NavyBlue,urlcolor=NavyBlue,citecolor=NavyBlue]{hyperref}

% Code listings
\usepackage{listings}
\lstset{
  basicstyle=\ttfamily\small,
  breaklines=true,
  frame=single,
  backgroundcolor=\color{gray!10},
  xleftmargin=2em,
  framexleftmargin=1.5em,
  numbers=none,
  showstringspaces=false,
  columns=flexible,
  literate={←}{$\leftarrow$}1 {→}{$\rightarrow$}1 {⟶}{$\longrightarrow$}1,
}
\lstdefinelanguage{json}{
  string=[s]{"}{"},
  comment=[l]{//},
  morecomment=[s]{/*}{*/},
  literate=
    *{:}{{{\color{NavyBlue}:}}}{1}
    {,}{{{\color{NavyBlue},}}}{1}
    {\{}{{{\color{NavyBlue}\{}}}{1}
    {\}}{{{\color{NavyBlue}\}}}}{1}
    {[}{{{\color{NavyBlue}[}}}{1}
    {]}{{{\color{NavyBlue}]}}}{1},
}
\lstdefinelanguage{toml}{
  comment=[l]{\#},
  string=[b]",
  morestring=[b]',
}
\lstdefinelanguage{protobuf}{
  morekeywords={service,rpc,returns,message,string,uint32,uint64,int64,bool,bytes,repeated,optional,oneof,enum,map},
  morecomment=[l]{//},
  morestring=[b]",
}

% Tables
\usepackage{longtable}
\usepackage{booktabs}
\usepackage{array}
\usepackage{tabularx}

% Graphics
\usepackage{graphicx}
\usepackage{float}

% Better lists
\usepackage{enumitem}

% Header/footer
\usepackage{fancyhdr}
\pagestyle{fancy}
\fancyhf{}
\fancyfoot[C]{\thepage}
\renewcommand{\headrulewidth}{0.4pt}

% Tighter spacing
\setlength{\parindent}{0pt}
\setlength{\parskip}{6pt}

% Convenience macros
\newcommand{\tcfs}{\texttt{tcfs}}
\newcommand{\tcfsd}{\texttt{tcfsd}}
\newcommand{\code}[1]{\texttt{#1}}
\newcommand{\filepath}[1]{\texttt{#1}}


\fancyhead[L]{\small\textsc{tcfs}}
\fancyhead[R]{\small Protocol Specification}

\title{tcfs Protocol Specification}
\author{tcfs / TummyCrypt}
\date{February 2026}

\begin{document}
\maketitle
\tableofcontents
\newpage

% ─────────────────────────────────────────────────────────────────────────────
\section{Stub File Formats}

\tcfs{} uses lightweight metadata to represent remote content locally. Physical
sync-root paths may use \code{.tc}/\code{.tcf} stub files, mounted VFS views
normally expose clean filenames backed by remote index entries, and platform
integrations such as FileProvider use native placeholders.

\subsection{\code{.tc} File Stub}

A physical \code{.tc} stub replaces a single file in sync-root/dehydrated paths
(e.g., \code{photo.jpg} becomes \code{photo.jpg.tc}). The stub is a sorted
key/value text file containing the metadata needed to reconstruct the original
file.

\begin{lstlisting}
version https://tummycrypt.io/tcfs/v1
chunks 23
compressed 0
fetched 0
oid blake3:4d7a214614ab2935c943f9e0ff69d22eadbb8f32b1258daaa5e2ca24d17e239
origin seaweedfs://filer.example.com/bucket/path/to/file
size 94371840
\end{lstlisting}

\begin{longtable}{lll}
  \toprule
  \textbf{Field} & \textbf{Type} & \textbf{Description} \\
  \midrule
  \endhead
  \code{version}    & string & Stub format URI \\
  \code{chunks}     & usize  & Number of chunks in the remote manifest \\
  \code{compressed} & bool   & Reserved for seekable-zstd; always \code{0} today (no compression in the data path) \\
  \code{fetched}    & bool   & \code{0} for dehydrated, \code{1} for hydrated \\
  \code{oid}        & string & Content OID, normally \code{blake3:\{hex\}} \\
  \code{origin}     & string & Remote origin URL \\
  \code{size}       & u64    & Original file size in bytes \\
  \bottomrule
\end{longtable}

\subsection{\code{.tcf} Folder Stub}

A \code{.tcf} stub is the physical directory-stub counterpart used by
dehydrated/compatibility paths. Mounted directory traversal should normally use
the remote index namespace instead of exposing raw \code{.tcf} entries.

% ─────────────────────────────────────────────────────────────────────────────
\section{Chunk Layout (S3/SeaweedFS)}

Content-addressed storage using BLAKE3 hashes as keys:

\begin{lstlisting}
{prefix}/
  chunks/
    {blake3_hash_1}     # Encrypted chunk (BLAKE3 of ciphertext; not compressed)
    {blake3_hash_2}
    ...
  manifests/
    {manifest_object_id}  # Manifest (plaintext JSON), path-indexed uploads
    {file_blake3_hash}    # Manifest for direct (non-indexed) uploads
  index/
    {rel_path}            # Index entry naming the current manifest
\end{lstlisting}

\subsection{Chunk Manifest Format}

The current regular-file manifest is \code{SyncManifest} version 2
(\code{crates/tcfs-sync/src/manifest.rs}), serialized as plaintext pretty
JSON:

\begin{lstlisting}[language=json]
{
  "version": 2,
  "file_hash": "<BLAKE3 hex of the plaintext file>",
  "file_size": 4194304,
  "chunks": ["<BLAKE3 hex of chunk ciphertext>", "..."],
  "vclock": { "...": 1 },
  "written_by": "<device id>",
  "written_at": 1760000000,
  "rel_path": "path/relative/to/root",
  "mode": 420,
  "mtime": [1760000000, 0],
  "encrypted_file_key": "<base64 wrapped file key>"
}
\end{lstlisting}

\code{chunks} is an ordered list of chunk keys only. Manifests do not record
chunk offsets or lengths, and nothing is compressed. The manifest itself is
not encrypted: the relative path, size and plaintext file hash are readable
by anyone who can read the bucket. \code{rel\_path} is optional and is
written as \code{null} when absent. \code{mode}, \code{mtime} and
\code{encrypted\_file\_key} are optional and omitted when absent.
\code{wrapped\_file\_keys} (per-device wraps) is omitted when empty, which
is the case under \code{wrap\_mode = master}.
Symlinks use a separate version-3 \code{SymlinkManifest}.

% ─────────────────────────────────────────────────────────────────────────────
\section{Chunking Algorithm}

\tcfs{} uses \textbf{FastCDC} (Fast Content-Defined Chunking) for splitting
files into variable-size chunks. The default profile (most files) is:

\begin{longtable}{ll}
  \toprule
  \textbf{Parameter} & \textbf{Value} \\
  \midrule
  Minimum chunk size & 2\,KiB \\
  Average chunk size & 4\,KiB \\
  Maximum chunk size & 16\,KiB \\
  Hash function      & BLAKE3 \\
  Compression        & none \\
  \bottomrule
\end{longtable}

Two larger profiles are selected by path (\code{ChunkSizes::for\_path} in
\code{crates/tcfs-chunks/src/fastcdc.rs}): 32/64/256\,KiB for the Git
index and for \code{.idx} and \code{.bin} files outside a Git pack
directory, and 1/4/16\,MiB for \code{.pack}, \code{.rev}, \code{.iso} and
\code{.img} files and, inside \code{.git/objects/pack/}, for \code{.idx}
files and extensionless \code{tmp\_pack\_*} files.

Content-defined chunking ensures that inserting or modifying bytes in a file
only affects the chunks near the edit point. With encryption on, each file
has its own random key and every chunk a random nonce, so identical plaintext
in two files does not produce identical stored chunks.

% ─────────────────────────────────────────────────────────────────────────────
\section{Hydration Flow}

When a remote-backed file is opened through the mounted VFS:

\begin{enumerate}
  \item The mount backend's \code{open()} handler intercepts the request.
  \item The index entry at \filepath{\{prefix\}/index/\{rel\_path\}} names the
        manifest, which is fetched from
        \filepath{\{prefix\}/manifests/\{manifest\_object\_id\}}.
  \item Chunks are fetched in parallel from
        \filepath{\{prefix\}/chunks/\{chunk\_hash\}}.
  \item Each chunk is decrypted (XChaCha20-Poly1305). There is no
        decompression step.
  \item BLAKE3 hash of the reassembled file is verified against \code{file\_hash}.
  \item The file is cached and served to the calling process via \code{read()}.
\end{enumerate}

% ─────────────────────────────────────────────────────────────────────────────
\section{State Tracking}

The sync engine maintains a local state cache (JSON file, migrating to RocksDB)
tracking per-file metadata:

\begin{itemize}
  \item Local path, remote S3 key, BLAKE3 hash, size
  \item Sync states: \code{synced}, \code{modified}, \code{pending\_upload},
        \code{pending\_download}, \code{conflict}
\end{itemize}

% ─────────────────────────────────────────────────────────────────────────────
\section{Wire Protocol (gRPC)}

The daemon (\tcfsd{}) exposes a gRPC service over a Unix domain socket:

\begin{lstlisting}[language=protobuf]
service TcfsDaemon {
  rpc Status(StatusRequest) returns (StatusResponse);
  rpc Mount(MountRequest) returns (MountResponse);
  rpc Unmount(UnmountRequest) returns (UnmountResponse);
  rpc Push(stream PushChunk) returns (stream PushProgress);
  rpc Pull(PullRequest) returns (stream PullProgress);
  rpc Hydrate(HydrateRequest) returns (stream HydrateProgress);
  rpc Unsync(UnsyncRequest) returns (UnsyncResponse);
  rpc SyncStatus(SyncStatusRequest) returns (SyncStatusResponse);
  rpc Watch(WatchRequest) returns (stream WatchEvent);
  rpc CredentialStatus(Empty) returns (CredentialStatusResponse);
  rpc ListRegisteredRoots(ListRegisteredRootsRequest)
      returns (ListRegisteredRootsResponse);
  rpc GetRegisteredRootStatus(GetRegisteredRootStatusRequest)
      returns (GetRegisteredRootStatusResponse);
  rpc ListConflicts(ListConflictsRequest)
      returns (ListConflictsResponse);
  rpc ResolveRegisteredRoot(ResolveRegisteredRootRequest)
      returns (ResolveRegisteredRootResponse);
  rpc ResolveConflict(ResolveConflictRequest)
      returns (ResolveConflictResponse);
}
\end{lstlisting}

For a non-primary root, \code{ListConflicts} carries a stable \code{root\_id}.
The separate \code{ResolveRegisteredRoot} method carries that ID, a checked
repository path, and one of the two repo-group Git keep-both modes (dry-run or
execute). The daemon resolves object prefix, state cache, and mutation policy
from trusted configuration and returns the selected route identity in the same
response. This scope does not provide named ordinary-file resolution or
\code{reconcile --root}. Legacy \code{ResolveConflict} has no root selector;
when named roots are enrolled, tcfsd fences rootless paths that enter them.

\subsection{Versioned Root Inventory (B0a)}

B0a adds a separate strict \code{sync.root\_registry} rather than changing the
meaning of the unversioned PR~\#551 conflict registry. Only explicit V1
descriptors appear in \code{ListRegisteredRoots} and
\code{GetRegisteredRootStatus}; legacy \code{sync.roots} entries are never
inferred, upgraded, or enumerated through these methods.

The fleet-stable \code{RootSpecV1} contains the version, root ID, remote prefix,
profile, generation, and a domain-separated BLAKE3 identity fingerprint. The
only V1 profiles are \code{git-raw-v1} and \code{agent-static-v1}. An optional
host-local \code{RootBindingV1} contains canonical local-root and state-cache
paths, lifecycle and resolution policy, and a separately domain-separated
binding fingerprint. This separation permits one fleet identity to bind to
unlike paths on different operating systems without making a host path part of
the fleet identity.

Both methods require a real authenticated daemon session with pull permission.
Rows are filtered by component-bounded registered-prefix authorization before
host filesystem inspection. A syntactically valid unknown ID and an
unauthorized ID produce the same generic not-found response.

Each row reports one availability state: unbound, unsupported platform,
missing local root, missing state, invalid binding, busy, or ready. A ready row
may include exclusive counts for the persisted states
\code{not\_synced}, \code{synced}, \code{active}, \code{locked}, and
\code{conflict}. These counts are an immutable snapshot of the registered
root's primary state-cache file, not live task telemetry. Inventory never
creates or modifies the sibling lock, repairs the primary state-cache file, or
recovers from its backup.

Every B0a response reports reconcile support as \code{NONE}. The methods do not
plan or execute reconcile, resolve conflicts, hydrate, unsync, enroll, or
mutate registry entries. No B0a method is exposed through MCP. The complete
source contract is recorded in
\filepath{docs/design/versioned-root-registry-status-b0a-2026-07-19.md}.

See \filepath{crates/tcfs-core/src/proto/tcfs.proto} for the full service
definition including all message types.

\end{document}
